% Source-derived chapter 11: Proof of \autoref{theorem:period-map-computation}
% Original source: canonical-representations-surface-groups
\chapter{Proof of \autoref{theorem:period-map-computation}}
\label{canonical-representations-surface-groups-chapter-11}
\source{canonical-representations-surface-groups}

\begin{remark}[Source section]
\label{canonical-representations-surface-groups-chapter-11-source}
\source{canonical-representations-surface-groups}
\source{canonical-representations-surface-groups}
This chapter reproduces the corresponding section of the retrieved
LaTeX source, including its definitions, statements, examples, and
proofs. It has not yet been translated into Lean.
\notready
\end{remark}

% The source section title was: Proof of \autoref{theorem:period-map-computation}
\label{appendix}
\source{canonical-representations-surface-groups}
We now explain the proof of \autoref{theorem:period-map-computation}, which loosely follows \cite[\S9 and \S10]{Voisin:hodgeTheory}. We retain notation as in \autoref{section:unitary-period-map}.
\subsection{Preparatory Lemmas}
Recall that the derivative of the period map $dP$ was a map 	$$dP : T_{\mathscr B} \to P^* T_{\on{Gr}(\rk F^1 \mathscr H, \rk \mathscr H)} \simeq (F^1
	\mathscr H)^\vee \otimes (\mathscr H/F^1 \mathscr H).$$
By adjointness, we obtain from $dP$ a map
\begin{align*}
	dP' : F^1 \mathscr H \to 
	(\mathscr H/F^1 \mathscr H) \otimes T_{\mathscr B}^\vee.
\end{align*}
Recall that we defined in \autoref{subsection:period-map-derivative} a map $\overline{\nabla}$ with the same source and target.
%At a point $b \in \mathscr B$, 
%using the canonical identification
%$T_{\on{Gr}(\rk F^1 \mathbb W, \rk \mathbb W), b}
%	\simeq \Hom(H^0(X_b, \mathscr E_b \otimes\Omega_{X_b/b}),H^1(X_b,
%	\mathscr E_b)),$
%this induces a map
%\begin{align}
%	\label{equation:derivative-period-at-b}
%	dP_b : T_{\mathscr B,b} \to T_{\on{Gr}(\rk F^1 \mathbb W, \rk \mathbb W), P(b)}
%	\simeq \Hom(F^1 \mathscr H, \mathscr H/F^1 \mathscr H)
%	H^0(X_b, \mathscr E_b \otimes\Omega_{X_b}),H^1(X_b,
%	\mathscr E_b)).
%\end{align}

\begin{lemma}
\source{canonical-representations-surface-groups}
\notready
	\label{lemma:connection-and-period-map}
\source{canonical-representations-surface-groups}
	We have an equality $\overline \nabla = dP'$.
\end{lemma}
\begin{proof}
	Let $W$ be a vector space and $K\subset W$ a subspace of rank $r$. It is shown in the course of the proof of \cite[Lemma 10.7]{Voisin:hodgeTheory}
	that for $[K] \in \on{Gr}(r,W)$, the identification $T_{\on{Gr}(r,W),[K]} \simeq
	\Hom(K,W/K)$ is given as follows.
	Let $\mathscr S \subset \mathscr{O}_{\on{Gr}(r,W)} \otimes W$ denote the universal
	subbundle.
	Given $\sigma \in K$, let $\widetilde{\sigma}$
	denote a choice of holomorphic section of $\mathscr S$
	defined on a neighborhood of $[K] \in \on{Gr}(r,W)$ with $\widetilde{\sigma}([K]) =
	\sigma$.
	Then, the map 
	\begin{align*}
		T_{\on{Gr}(r, W), [K]} &\to \Hom(K, W/K) \\
		u &\mapsto \left( \sigma \mapsto \frac{\partial }{\partial
		u} (\widetilde{\sigma}) \bmod K \right)
	\end{align*}
	is well-defined and independent of the choice of lift $\widetilde{\sigma}$.

	Choosing $b \in \mathscr B$, the derivative of the period map $dP_b'$ is then given as follows:
	send 
	$\sigma \in F^1 \mathscr H_b$ to the function which sends
	$v \in T_{\mathscr B,b}$ to $\frac{\partial }{\partial v} \widetilde{\sigma} \bmod F^1
	\mathscr H_b$,
	for $\widetilde{\sigma}$ a local holomorphic lift of $\sigma$ to $F^1
	\mathscr H$. Here $\frac{\partial }{\partial v}$ makes sense as we have
	chosen a flat trivialization of $\mathscr{H}$, but we may equivalently
	write this map as $$\sigma \mapsto (v \mapsto \nabla_{GM}(\widetilde{\sigma})(v)
	\bmod F^1\mathscr{H}_b),$$
	or equivalently 
	\begin{equation*}
	\sigma \mapsto (v \mapsto
	\overline{\nabla}(\widetilde{\sigma})(v)).\qedhere
\end{equation*}
%
%	We wish to show $dP'(\sigma)(v) =
%	\overline{\nabla}(\sigma)(v),$ and so it suffices to show that
%	$\overline{\nabla}$
%	is also given by sending a section $\sigma$ to $\frac{\partial
%	}{\partial v} \widetilde{\sigma} \bmod F^1 \mathscr H$, for $\widetilde{\sigma}$ a
%	holomorphic lift of $\sigma$.  
%	Because our connection $\nabla$ constructed as $\on{id} \otimes d$, it is
%	given by sending $\widetilde{\sigma}$ to the function sending $v$ to
%	$\frac{\partial }{\partial v} \widetilde{\sigma}$. 
%	Here, we take
%	$\widetilde{\sigma}$ to be a holomorphic section of $W$ restricting to $\sigma$
%	at $b$. Since this is a holomorphic lift of $\sigma$, we precisely recover $dP'(\sigma)(v) =
%	\overline{\nabla}(\sigma)(v),$ as desired.
%	Now, to compute $\overline{\nabla}$ over a point $b$, we may choose a local
%	analytic neighborhood $U$ of $b$ trivialization $\mu: \mathbb W|_U \simeq
%	\underline{\mathbb C}_U^{\oplus \rk \mathbb W}$.
%	Upon taking this trivialization, the connection $\nabla = 1 \otimes d$
%	on $\mathbb W$ becomes simply the usual exterior derivative on 
%	$\underline{\mathbb C}_U^{\oplus \rk \mathbb W} \otimes_{\mathbb C}
%	\mathscr O_U$. That is, 
%	\begin{equation}
%		\label{equation:}
%		\begin{tikzcd} 
%			\mathbb W_U \otimes_{\mathbb C} \mathscr O_B \ar
%			{r}{\nabla = \on{id}
%			\otimes d} \ar
%			{d} & \mathbb W_U \otimes_{\mathbb C} \Omega_B^1 \ar
%			{d}{\mu \otimes \id} \\
%			\underline{\mathbb C}^{\oplus \rk \mathbb W}_U
%			\otimes_{\mathbb C} \mathscr O_B  \ar {r}{\on{id} \otimes
%			d}{\mu \otimes \on{id}} & 
%			\underline{\mathbb C}^{\oplus \rk \mathbb W}_U
%			\otimes_{\mathbb C} \Omega_B^1
%	\end{tikzcd}\end{equation}
%	commutes.
%	Therefore, the connection $\nabla$ is carried over to the directional
%	derivative on the local trivialization. Above, working on the local
%	trivialization of the Grassmannian, we identified this directional
%	derivative with the tangent space to the Grassmannian.
%	Pulling back along the period map $P$ identifies this directional
%	derivative with the derivative of the period map, as desired.
\end{proof}

After introducing some notation, we next give an explicit computation of the Gauss-Manin connection $\nabla_{GM}$ on $\mathscr{H}$. 

\begin{notation}
\source{canonical-representations-surface-groups}
\notready
	\label{notation:}
\source{canonical-representations-surface-groups}
	Recall from \autoref{notation:period-map}, that $(\mathscr E, \nabla)$ denotes the Deligne canonical extension of
$(\mathbb V \otimes \mathscr O_{\mathscr C^\circ}, \on{id} \otimes d)$ to
$\mathscr C$.
	Let $\mathcal{A}^{i,j}_{\log \mathscr{D}}(\mathscr{E})$ be the sheaf of $C^\infty$ logarithmic $(i,j)$-forms valued in $\mathscr{E}$, and let $\mathcal{A}^n_{\log \mathscr{D}}(\mathscr{E})=\oplus_{i+j=n} \mathcal{A}^{i,j}_{\log \mathscr{D}}(\mathscr{E})$. 
	Then, $(\mathcal{A}^{i,\bullet}_{\log\mathscr{D}}(\mathscr{E}),
	\overline{\partial})$ is the Dolbeault complex of $\mathscr{E}\otimes
	\Omega^i_{\mathscr{C}}(\log \mathscr D)$. The complex
	$(\mathcal{A}^\bullet_{\log \mathscr{D}}(\mathscr{E}),
	\nabla+\overline{\partial})$ is the de Rham resolution of the unitary
	local system $\mathbb{V}$. We similarly define the relative variants  $\mathcal{A}^{i,j}_{\log \mathscr{D}, \mathscr{C}/\mathscr{B}}(\mathscr{E})$, $\mathcal{A}^{n}_{\log \mathscr{D}, \mathscr{C}/\mathscr{B}}(\mathscr{E})$.
\end{notation}
\begin{notation}
\source{canonical-representations-surface-groups}
\notready
	\label{notation:}
\source{canonical-representations-surface-groups}
	For $v$ a $C^\infty$ vector field on an open $U \subset \mathscr C$, tangent to $\mathscr{D}$, and
$w$ a $C^\infty$ section of $\mathcal{A}^i_{\log \mathscr{D}}(\mathscr{E})$ on $U$, the interior product
$\on{int}(v)(w) \in H^0(U, \mathcal{A}^{i-1}_{\log \mathscr D}(\mathscr E))$ is
defined by
$$\on{int}(v)(w)(X_1, \ldots, X_{i+j-1}) = w(v, X_1, \ldots, X_{i+j-1})$$ for any
$C^\infty$ vector fields $X_1, \ldots, X_{i+j-1}$ on $U$ tangent to $\mathscr{D}$.
\end{notation}

\begin{lemma}
\source{canonical-representations-surface-groups}
\notready
	\label{lemma:connection-to-interior}
\source{canonical-representations-surface-groups}
	Let $\sigma\in \mathcal A^1_{\log \mathscr{D}}(\mathscr{E})$ be a
	differential form such that for each $b'\in \mathscr{B}$,
	$\sigma|_{\mathscr{C}_{b'}}$ is closed, so that $\sigma$ represents a
	section $[\sigma]$ of $\mathscr{H}$.
Fix $b\in \mathscr{B}$ and a tangent vector $u\in T_{\mathscr{B}, b}$, and let
$v$ be a $C^\infty$ section of $T_\mathscr{C}|_{\mathscr{C}_b}$, tangent to
$\mathscr{D}$ with $u = \pi_*(v).$
	We have
	\begin{align*}
		\nabla_{GM}([\sigma])|_b(u)= [(\on{int}(v)((\nabla+\overline{\partial})
	(\sigma)|_{\mathscr C_b})]\in \mathscr{H}_b.
	\end{align*}
\end{lemma}
\begin{proof}
	This proof essentially follows \cite[Proposition
	9.14]{Voisin:hodgeTheory}.	
	After possibly shrinking $\mathscr B$, which will not alter the statement of this
	lemma, we may write $\mathscr C \simeq
	\mathscr C_b \times \mathscr B$ (as $C^\infty$ manifolds), by Ehresmann's theorem. 
	Let $t_1, \ldots, t_m$ be $C^\infty$  functions on $C_b \times
	\mathscr B$ pulled back from a system of local coordinates for $\mathscr B$.
	After possibly shrinking $\mathscr B$, we can write
	$\sigma$ as
	\begin{align*}
		\sigma = \Phi + \sum_{i=1}^m dt_i \otimes \psi_i, 
	\end{align*}
	where $\Phi$ a section of $\mathscr{A}^1_{\log \mathscr{D}}(\mathscr{E})$ independent of
	$dt_i$, 
	each
	$\psi_i$ is a $C^\infty$ section of $\mathscr{E}$,
	and $\Phi|_{\mathscr C_b} = \sigma|_{\mathscr C_b}$. 
	Since $\sigma$ is fiber-wise closed, applying $\nabla+\overline{\partial}$ gives
		\begin{align*}
		(\nabla+\overline{\partial}) \sigma = \sum_i dt_i
		\on{int}(\partial/\partial t_i)((\nabla+\overline{\partial})\Phi) - \sum_i dt_i \wedge (\nabla+\overline{\partial})
		\psi_i 
	\end{align*}
	and hence
	\begin{align*}
		\on{int}(\partial/\partial t_i) ((\nabla+\overline{\partial})\sigma)|_{\mathscr{C}_b} =
	\on{int}(\partial/\partial t_i)((\nabla+\overline{\partial})\Phi|_{\mathscr{C}_b}) - (\nabla+\overline{\partial})
		\psi_i|_{\mathscr{C}_b}. 
	\end{align*}
	
	Since $\Phi|_{\mathscr C_b} = \sigma|_{\mathscr C_b}$, 
	$$\on{int}(\partial/\partial
	t_i)((\nabla+\overline{\partial})\Phi|_{\mathscr{C}_b}) =
	\nabla_{GM}([\sigma])|_b(\pi_*(\partial/\partial t_i)),$$ by definition of the Gauss-Manin connection.	
	But as $(\nabla+\overline{\partial})
		\psi_i|_{\mathscr{C}_b}$ is exact, 
		$\on{int}(\partial/\partial t_i) ((\nabla+\overline{\partial})\sigma)|_{\mathscr{C}_b}$ and $\on{int}(\partial/\partial t_i)((\nabla+\overline{\partial})\Phi|_{\mathscr{C}_b})$ represent the same cohomology class. Thus 
\begin{align*}
	\mathscr{H}_b \ni [(\on{int}(\partial/\partial t_i)((\nabla+\overline{\partial})(\sigma)|_{\mathscr C_b})] =
		[(\on{int}(\partial/\partial t_i)((\nabla+\overline{\partial})(\Phi)|_{\mathscr C_b})] =
		\nabla_{GM}([\sigma])|_b(\pi_*(\partial/\partial t_i)).
\end{align*}

 Since
	$\sigma|_{\mathscr{C}_b}$ is closed, 
	the same result holds after
	replacing $\partial/\partial t_i$ with any vector field $v$ satisfying
	$\pi_*v=\pi_*(\partial/\partial t_i)$; 	
	as the $\pi_*(\partial/\partial t_i)$ form a
	basis for $T_{\mathscr{B},b}$, the proof is complete.
		\end{proof}

For the next result, we need to recall the Kodaira-Spencer map.
\begin{definition}
\source{canonical-representations-surface-groups}
\notready
	\label{definition:kodaira-spencer}
\source{canonical-representations-surface-groups}
	With notation as in \autoref{notation:period-map},
we have a short exact sequence
\begin{equation}
	\label{equation:tangent-exact-sequence}
\source{canonical-representations-surface-groups}
	\begin{tikzcd}
		0 \ar {r} & T_{\mathscr C_b}(-\log \mathscr D) \ar {r} & T_{\mathscr C}(-
		\log \mathscr D)|_{\mathscr C_b} \ar {r}
	& \pi^* T_{\mathscr B}|_{\mathscr C_b} \ar {r} & 0 
\end{tikzcd}\end{equation}
where $T_{\mathscr{C}}(-\log\mathscr{D})$ is the dual of $\Omega^1_{\mathscr{C}}(\log \mathscr{D})$.
The {\em Kodaira-Spencer} map $$\kappa: T_{\mathscr{B}, b}\to H^1(T_{\mathscr{C}_b}(-\log\mathscr{D}))$$ sends 
a vector $v \in T_{\mathscr B,b} \simeq H^0(\mathscr C_b, \pi^* T_{\mathscr B}|_{\mathscr C_b})$ to its image in 
$H^1(T_{\mathscr{C}_b}(-\log\mathscr{D}))$ under the connecting homomorphism in the long exact sequence in cohomology arising from \eqref{equation:tangent-exact-sequence}. 
\end{definition}


\begin{remark}
\source{canonical-representations-surface-groups}
\notready
	\label{remark:kodaira-spencer}
\source{canonical-representations-surface-groups}
	Geometrically, the Kodaira-Spencer map sends a deformation of $(\mathscr C_b, \mathscr D_b)$ to the corresponding
cohomology class in $H^1(T_{\mathscr C_b}(-\mathscr D))$ (identifying $T_{\mathscr C_b}(-\mathscr D)$ with $T_{\mathscr{C}_b}(-\log\mathscr{D})$).
For an explanation of why the boundary map coincides with this description in
the case $\mathscr D = \emptyset$, see, \cite[\S9.1.2]{Voisin:hodgeTheory}.
In the case $\mathscr D \neq \emptyset$, an analogous argument goes through.
\end{remark}

\begin{proposition}
\source{canonical-representations-surface-groups}
\notready
	\label{proposition:period-map}
\source{canonical-representations-surface-groups}
	Keep notation as in \autoref{notation:period-map}.
	For $b \in \mathscr B$, set $C=\mathscr{C}_b$, $D=\mathscr{D}_b$, and $E=\mathscr{E}|_C$. We have a commuting diagram
	\begin{equation}
		\label{equation:period-to-cup}
\source{canonical-representations-surface-groups}
		\begin{tikzcd} 
			H^0( C, E \otimes \Omega^1_{ C}(D)) \otimes
			T_{\mathscr B,b} \ar {r}{dP_b''} \ar {d}{1 \otimes \kappa} & H^1( C,  E) \\
			H^0( C,  E \otimes \Omega^1_{ C}( D)) \otimes
			H^1( C, T_{ C}(- D)) \ar {r}{\cup} & H^1( C,  E
				\otimes \Omega^1_{ C}( D) \otimes T_{ C}(- D)) \ar{u}{\alpha}
	\end{tikzcd}\end{equation}
	where $dP_b''$ is adjoint to the derivative of the period map $dP_b$, $\kappa$ is the Kodaira-Spencer map, $\alpha$ is the map induced by the pairing 
	$\Omega^1_{ C}( D) \otimes T_{ C}(- D) \to \mathscr O_{ C}$
	and $\cup$ is the cup product.
\end{proposition}
\begin{proof}
Let $\sigma\otimes u$ be an element of $H^0( C, E \otimes \Omega^1_{ C}(D)) \otimes
			T_{\mathscr B,b}$. By e.g.~\cite[Proposition
			9.22]{Voisin:hodgeTheory} (applied to the
			$\overline\partial$-Laplacian---here unitarity is used
		to define the adjoint of $\overline\partial$ and consequently
	the Laplacian), we may choose a lift $\widetilde\sigma$ of $\sigma$ to a $C^\infty$ section of $\mathscr{E}\otimes \Omega^1_{\mathscr{C}}(\log \mathscr{D})$, which is holomorphic on fibers.
	Combining \autoref{lemma:connection-and-period-map} and
	\autoref{lemma:connection-to-interior}, we find that $dP''(\sigma
	\otimes u)$
	can be expressed as $[((\on{int}(v)(\nabla+\overline{\partial})\widetilde{\sigma})|_{\mathscr C_b})^{0,1}]$ for
$v$ a $C^\infty$ vector field of type $(1,0)$ tangent to $\mathscr{D}$ and lifting $u$, as in \autoref{lemma:connection-to-interior}. 
As $\mathscr{C}$ is a relative curve, for degree reasons we have
$$[((\on{int}(v)(\nabla+\overline{\partial})\widetilde{\sigma}|_{\mathscr C_b}))^{0,1}]=[\on{int}(v)\overline{\partial}\widetilde{\sigma}|_{\mathscr C_b}].$$
Direct computation gives
\begin{align*}
	\overline{\partial} (\on{int}(v) (\widetilde{\sigma}|_{\mathscr{C}_b})) = 
	- \on{int}(v)\overline{\partial}\widetilde{\sigma}|_{\mathscr C_b}
	+ \on{int}(\overline{\partial} v)(\widetilde{\sigma}|_{\mathscr{C}_b}).
\end{align*}
Using the vanishing of 
the Dolbeault cohomology class $[\overline{\partial} (\on{int}(v)
(\widetilde{\sigma}))]$
we obtain 
\begin{align*}
	[\on{int}(v)(\overline{\partial}\widetilde{\sigma}|_{\mathscr C_b})]
	=
	[\on{int}(\overline{\partial}v)(\widetilde{\sigma}|_{\mathscr C_b})].
\end{align*}
Combining the above yields that $dP''_b(\sigma\otimes u) = 
[\on{int}(\overline{\partial}v)(\widetilde{\sigma}|_{\mathscr C_b})]$.
By definition of the Kodaira Spencer map 
%(and the construction of the boundary map in cohomology)
$\overline{\partial}v$ is identified with $\kappa(u)$.
So
\begin{align*}
dP''_b(\sigma \otimes u) =
[\on{int}(\overline{\partial}v)(\widetilde{\sigma}|_{\mathscr C_b})]
=
[\alpha\left( \overline{\partial}v \cup \widetilde{\sigma}|_{\mathscr C_b}\right)]
=
[\alpha\left( \overline{\partial}v\cup \sigma \right)]
=
[\alpha\left( \kappa(u) \cup \sigma \right)]
\end{align*}
verifying the commutativity of
\eqref{equation:period-to-cup} as desired.
\end{proof}

\begin{lemma}
\source{canonical-representations-surface-groups}
\notready
	\label{lemma:multiplication-map}
\source{canonical-representations-surface-groups}
	Suppose $E$ is a vector bundle on a smooth curve $C$, and $D \subset C$
	is a divisor.
	Let $H^1(C, T_C(-D)) \otimes H^0(C, E \otimes \omega_C(D)) \to H^1(C, E)$
	denote the map corresponding to the Yoneda cup product composition $\alpha \circ \cup$ 
	from \autoref{proposition:period-map},
	induced by the natural pairing $E \otimes \omega_C(D) \otimes T_C(-D) \to E$.
This is adjoint to a map $\upsilon: H^0(C, E \otimes \omega_C(D)) \otimes
	H^1(C, E)^\vee \to H^1(C, T_C(-D))^\vee$.
	We have a commutative diagram
	\begin{equation}
		\label{equation:serre-dual-diagram}
\source{canonical-representations-surface-groups}
		\begin{tikzcd} 
			H^0(C, E \otimes \omega_C(D)) \otimes H^0(C, E^\vee \otimes
			\omega_C)  \ar {r}{\mu}\ar{d}{\eta}  & H^0(C, \omega_C^{\otimes 2}(D))\ar{d}{\zeta}
\\
H^0(C, E \otimes \omega_C(D)) \otimes H^1(C, E)^\vee \ar
			{r}{\upsilon} & H^1(C, T_C(-D))^\vee ,
				\end{tikzcd}\end{equation}
	where the two vertical maps are induced by Serre duality, and the 
	map $\mu$ is the map on global sections induced by $$(E
	\otimes \omega_C(D)) \times (E^\vee \otimes \omega_C) \to (E \otimes E^\vee)
	\otimes (\omega_C(D) \otimes \omega_C) \to \omega_C^{\otimes 2}(D)$$
	where the first map is the tensor product and the second is obtained from the trace pairing $E\otimes E^\vee\to \mathscr{O}_C$.
\end{lemma}
\begin{proof}
	Recall that for $F$ a vector bundle on $C$, the Serre duality pairing
	between $H^0(C, F)$ and $H^1(C, F^\vee\otimes \omega_C)$ is obtained
	as the composition $H^0(C, F) \otimes H^1(C, F^\vee\otimes \omega_C) \to
	H^1(C, F^\vee\otimes F \otimes \omega_C) \to
	H^1(C, \omega_C) \to \mathbb C$, where
	the first map is induced by cup product, the second is induced by the
	pairing
	$F \otimes F^\vee \to \mathscr O$,
	and the third map is the trace
	map $\on{tr}: H^1(C, \omega_C) \simeq \mathbb C$
	\cite[30.3.15]{vakil:foundations-of-algebraic-geometry-2017}.

	Using the above description, we now check commutativity of
	\eqref{equation:serre-dual-diagram}.
	We use $\cup$ for the usual cup product, and $\overline{\cup}$ to denote
	the composition of the cup product and the pairing $E \otimes E^\vee \to
	\mathscr O$. 
		Choose $\alpha \in H^0(C, E \otimes \omega_C(D)), \beta \in
	H^0(C, E^\vee \otimes \omega_C)$; we will denote by $s$ a general element of $H^1(C, E)$ and by $t$ a general element of $H^1(C,
	T_C(-D))$.
	On the one hand $\mu(\alpha \otimes \beta) = \alpha \overline{\cup}\beta$, and so $\zeta(\mu(\alpha\otimes \beta))$ is given by $$t\mapsto \on{tr}((\alpha\overline{\cup}\beta)\cup t).$$
	On the other hand, using the above description of Serre duality,
	\begin{align*}
		\upsilon(\eta(\alpha \otimes \beta))
		&= \upsilon(\alpha \otimes (s \mapsto \on{tr}(\beta
		\overline{\cup}s)))
		\\
		&= \left(t \mapsto \on{tr}(\beta \overline{\cup} (\alpha \cup
			t))\right) \\
			&= \left( t \mapsto \on{tr}( (\beta \overline{\cup} \alpha) \cup
			t )\right)  
	\end{align*}
	and so we obtain the desired commutativity (using that $\alpha \overline{\cup} \beta = \beta \overline{\cup} \alpha$).
%%The following is a more diagrammatic proof, but it is much longer, and probably
%%	more confusing.
%	A diagram chase gives that the commutativity of \eqref{equation:serre-dual-diagram}
%	is equivalent to commutativity of
%		\begin{equation}
%	\label{equation:cup-product-diagram}
%	\adjustbox{scale=.8,center}{
%		\begin{tikzcd} 
%			H^1(C, T_C(-D)) \otimes H^0(C, E \otimes \omega_C(D)) \otimes
%			H^0(C, E^\vee \otimes \omega_C) \ar {d}{\beta} \ar
%			{r}{\gamma} & 
%			H^1(T_C(-D)) \otimes H^0(C, \omega_C^{\otimes 2}(D))
%			\ar {d}{\delta} \\
%			H^1(E)\otimes H^0(C, E^\vee \otimes \omega_C) 
%			 \ar {r}{\varepsilon} & H^1(\omega_C), 
%		\end{tikzcd}
%	}
%	\end{equation}
%	where the maps in this diagram are defined as follows:
%	$\beta$ is induced by the cup product on the first two factors,
%	$\gamma$ is induced by the multiplication map $\mu$,
%	$\delta$ is induced by the cup product, and $\varepsilon$ is induced
%	by the cup product and the trace pairing $E \otimes E^\vee \otimes \omega_C \to \omega_C$.
%
%	We now verify commutativity of
%	\eqref{equation:cup-product-diagram}.
%	Diagram \eqref{equation:cup-product-diagram}
%	 is induced by cup products arising from the
%	commutative diagram of sheaves
%	\begin{equation}
%		\label{equation:}
%		\begin{tikzcd} 
%			T_C(-D) \otimes (E \otimes \omega_C(D)) \otimes (E^\vee \otimes
%		\omega_C) \ar {r} \ar {d} & T_C(-D) \otimes \omega_C^{\otimes
%		2}(D) \ar {d} \\
%		E \otimes (E^\vee \otimes \omega_C) \ar {r} & \omega_C.
%	\end{tikzcd}\end{equation}
%	Here, the maps of sheaves are the natural ones, arising from pairing dual vector bundles together. Commutativity of \eqref{equation:cup-product-diagram} follows from the associativity of the cup product.
%	We conclude by verifying the equivalence of commutativity of
%	\eqref{equation:serre-dual-diagram} and
%	\eqref{equation:cup-product-diagram}.
%	First, note that by pairing with an element of $H^1(C, T_C(-D))$,
%	commutativity of \eqref{equation:serre-dual-diagram} is equivalent to
%	commutativity of an adjoint diagram
%	\begin{equation}
%		\label{equation:}
%		\begin{tikzcd} 
%			H^1(C, T_C(-D)) \otimes H^0(C, E \otimes \omega_C(D)) \otimes
%			H^0(C, E^\vee \otimes \omega_C) \ar {r}{\mu'}
%			\ar{d}{\eta'} & H^1(C, T_C(-D))
%			\otimes H^0(C, \omega_C^{\otimes 2}(D)) \ar{d}{\zeta'} \\
%			H^1(C, T_C(-D)) \otimes H^0(C, E \otimes \omega_C(D)) \otimes
%			H^1(C, E)^\vee \ar {r}{\upsilon'} & H^1(C,
%			\omega_C(D)).
%	\end{tikzcd}\end{equation}
%	where we have applied the adjunction and replaced the copy of $\mathbb
%	C$ in the upper right hand corner by $H^1(C, \omega_C)$ by composing
%	with $\on{tr}^{-1}$.
%	It is therefore enough to verify $\zeta' \circ \mu' = \delta \circ \gamma$ and
%	$\upsilon' \circ \eta'=\varepsilon \circ \beta$.
%	In fact, by construction, $\mu' = \gamma$.
%	Also, $\delta = \zeta'$ since the Serre duality map $\zeta$ is given as
%	the composition of $\zeta'$ with the trace map.
%	These together verify $\zeta' \circ \mu' = \delta \circ \gamma$.
%
%	We next show 
%	$\upsilon' \circ \eta'=\varepsilon \circ \beta$.
%	To verify this, we consider the diagram
%\begin{tikzpicture}[baseline= (a).base]
%\node[scale=.7] (a) at (0,0){
%				\begin{tikzcd}
%					H^1(C, T_C(-D)) \otimes H^0(E \otimes \omega_C(D))
%			\otimes H^0(E^\vee \otimes \omega_C) \ar{dd}{\eta'}
%			\ar{rr}{\beta} & &  
%			H^1(C, E) \otimes H^0(C, E \otimes \omega_C(D))
%			\ar{dd}{\varepsilon}
%			\\
%			& H^1(C, E) \otimes H^1(C,E)^\vee \ar{ur}{\nu}
%			\ar{rd}{\sigma}& \\
%			H^1(C, T_C(-D)) \otimes H^0(C, E \otimes \omega_C(D)) \otimes
%			H^1(C, E)^\vee \ar {ur}{\tau} \ar{rr}{\upsilon'} && 
%			H^1(C, \omega_C),
%			\\
%					 \end{tikzcd}   
%};
%\end{tikzpicture}
%where $\tau$ is given the Yoneda cup product by multiplication and $\nu$ is given by Serre duality.
%We will separately show $\beta = \nu \circ \tau \circ \eta'$ and $\sigma
%\circ \tau = \upsilon'$, and $\varepsilon \circ \nu = \sigma$.
%
%First we check $\beta = \nu \circ \tau \circ \eta'$
%Indeed, the map $\beta$ is given by the Yoneda cup product on the first two
%factors, while $\tau$ is similarly given by the Yoneda cup product, and $\eta'$
%and $\nu$ are given both by applying Serre duality to the third factor.
%Since applying Serre duality to the third factor commutes with Yoneda cup
%product on the first two factors, we obtain 
%$\beta = \nu \circ \tau \circ \eta'$.
%
%Next, 
%note that $\sigma \circ \tau = \upsilon'$ by construction of $\upsilon'$.
%To conclude, it remains to verify 
%$\varepsilon \circ \nu = \sigma$.
%But recall that $\sigma$ was by definition the composition of the natural pairing
%into $\mathbb C$ with the inverse of the trace map $\on{tr}$.
%Since $\varepsilon$ is the Yoneda cup product, $\on{tr} \circ \varepsilon$
%agrees with the Serre duality pairing given by $(\on{tr} \circ \sigma)
%\nu^{-1}$,
%and hence $\varepsilon = \sigma \circ \nu^{-1}$ so $\varepsilon \circ \nu =
%\sigma$.
	\end{proof}

\subsection{Proof of \autoref{theorem:period-map-computation}}
\begin{proof}
\label{subsection:proof-period-map}
\source{canonical-representations-surface-groups}
	The final statement claiming adjointness follows from
	\autoref{lemma:connection-and-period-map}.
	The map $\mu$ in \autoref{lemma:multiplication-map}, 
	is another name for the composition	
	$\on{tr} \circ \otimes$ in \eqref{equation:multiplication-to-period},
	and by \autoref{lemma:multiplication-map},
	$\mu$ is identified with the map labeled $\upsilon$ there, upon applying
	Serre duality.
	Moreover, $\upsilon$ is by definition adjoint to the composition $\alpha \circ \cup$
	in \autoref{proposition:period-map}.
	The result now follows by identifying the Kodaira-Spencer map $\kappa$
	of \autoref{proposition:period-map} as Serre dual to the pullback map
	$c_b^*$ in the statement of \autoref{theorem:period-map-computation};
	this identification follows from \autoref{remark:kodaira-spencer}.
\end{proof}

\bibliographystyle{alpha}
\bibliography{bibliography-mcg-hodge-theory}
